%! suppress = EscapeHashOutsideCommand
\section{Basics}

\bd[Command]
A \textbf{command} is a directive, in the form of lines of text, to a computer program to perform a specific task.
\ed

Commands can be followed by one or more ``options'' that modify their behavior and, further, by one or more
``arguments'', the items upon which the command acts.

\bd[Option / Flag]
An \textbf{option} or \textbf{flag} modifies the operation of a command; the effect is determined by the command's
program.
\ed

Options follow the command name separated by spaces and started by either a single dash \code{-} and the single
character version of the option or a double dash \code{--} and the long version of the option. Both single character
and long format are identical options and do the exact same thing. Also, many commands allow multiple short options to
be strung together.

\bd[Argument]
An \textbf{argument} is an item of information provided to a command when it is started.
\ed

Arguments follow the options (if any) or the command (if no options) separated by spaces. A command can have many
arguments that identify sources or destinations of information, or that alter the operation of the command. \v

So, most commands look kind of like this:
\begin{bash}
# common format of commands
command -options arguments
\end{bash}

A user can provide commands to a ``command-line interface''

\bd[Command-Line Interface (CLI)]
A \textbf{command-line interface (CLI)]} is a text-based user interface (UI) used to accept commands.
\ed

CLI runs programs, manages computer files and interacts with the computer through the ``command-line interpreter''.

\bd[Command-Line Interpreter]
A \textbf{command-line interpreter} is a computer program that exposes an operating system's services to a human user
or other programs, by receiving keyboard commands through a CLI and passing them to the operating system to carry out.
\ed

When we speak of the command line interpreter, we are really referring to the shell.

\bd[Shell]
The \textbf{shell} is a command-line interpreter acting as the outermost layer around the operating system.
\ed

The name ``shell'' for a command-line interpreter and the concept of making the shell a user program outside of the
operating system kernel were introduced in a Unix's precursor called Multics. \v

Since shells are just computer programs, it makes sense that there are a lot of them. The very fist shell of Unix was
the so called ``Thompson shell''.

\bd[Thompson Shell]
The \textbf{Thompson shell} was the first Unix shell, introduced in the first version of Unix in 1971, and was written
by Ken Thompson.
\ed

Thompson shell was a simple command line interpreter, not designed for scripting, but nonetheless introduced several
innovative features to the command-line interface and led to the development of the later Unix shells. The successor of
Thompson shell was the so called ``Bourne shell''.

\bd[Bourne Shell]
The \textbf{Bourne shell} is a Unix shell, released in 1979 in the Version 7 of Unix, developed by Stephen Bourne at
Bell Labs and acting as a replacement for the Thompson shell.
\ed

Bourne shell was used as an interactive command interpreter and it was also intended as a scripting language and
contains most of the features that are commonly considered to produce structured programs. Unix-like systems continue
to have a Bourne shell, or a symbolic link or hard link to a compatible shell—even when other shells are used by
most users. \v

Coming to today, almost all Unix systems supply a shell program called ``bash'' which has been used as the default login
shell for most of them.

\bd[Bash]
\textbf{Bash} is a Unix shell and command language, released in 1989, and written by Brian Fox for the GNU
Project\footnote{The GNU Project is a free software, mass collaboration project announced by Richard Stallman on
September 27, 1983. Its goal is to give computer users freedom and control in their use of their computers and computing
devices by collaboratively developing and publishing software that gives everyone the rights to freely run the software,
copy and distribute it,study it, and modify it.} as a free software replacement for the Bourne shell.
\ed

The name bash is an acronym for ``bourne-again shell'', a reference to the fact that bash is an enhanced replacement
for the original Unix shell program; Bourne. \v

Bash was also the default shell in versions of Apple macOS until it changed to ``Z shell'' or ``zsh'' (although Bash
remains available as an alternative shell).

\bd[Z Shell (zsh)]
The \textbf{Z shell (zsh)} is a Unix shell, developed in 1990, by Paul Falstad.
\ed

The name zsh derives from the name of Yale professor Zhong Shao. Paul Falstad regarded Shao's login-id, ``zsh'', as a
good name for a shell. Zsh is an extended Bourne shell with many improvements, including some features of Bash. \v

Moving on, in order to get access to and interact with a shell we need another program called ``terminal''.

\bd[Terminal Emulator / Terminal Application / Terminal]
A \textbf{terminal emulator}, or \textbf{terminal application}, or simply \textbf{terminal} is a graphical user
interface (GUI) computer program that gives us access to the shell.
\ed

A huge number of terminals are available, but they all basically do the same thing: give us access to the shell. The
most famous ones include ``Linux console'' for Linux, ``MacTerminal'' for Mac and ``iTerm2'' which is an open-source
terminal specifically for macOS. \v

Once a terminal is launched the first thing one sees is the so called ``shell prompt''.

\bd[Shell Prompt]
A \textbf{shell prompt} appears in the terminal whenever the shell is ready to accept input.
\ed

The default appearance of a shell prompt (although it might vary depending on the distribution), will typically include
your \code{username@machinename}, followed by the current working directory (more about that in a little bit) and a
dollar sign (\$):
\begin{bash}
[username@machinename current_directory]$\$$
\end{bash}

If the last character of the prompt is a hash mark (\#) rather than a dollar sign, the terminal session has superuser
privileges (more about that later). \v

Nowadays, the shell prompt's appearance can be configured to anyone's likings, so more often than not one can see many
different variations of a shell prompt. \v

Once we give a command to the terminal, this command is saved to the ``command line history''.

\bd[Command Line History]
\textbf{Command line history} is a feature in many operating shells that allows the user to recall, edit and rerun
previous commands.
\ed

Command line history was added to Unix by Bill Joy in 1978. It quickly became popular because it made the shell fast
and easy to use. Command line history has since become a standard feature in almost every shell. \v

History addresses two important scenarios:
\bit
\item Executing the same command or a short sequence of commands over and over.
\item Correcting mistakes or rerunning a command with only a small modification.
\eit

To access command line history we press the up and down arrows. If we try the left and right arrows, we’ll see that we
can position the cursor anywhere on the command line. This makes editing commands easy.

\section{Navigation}

\bd[File]
A \textbf{file} is a computer resource for recording data in a computer storage device.
\ed

Just as words can be written to paper, so can data be written to a file. Each file is distinguished from any other file
by its ``filename''.

\bd[Filename]
\textbf{Filename} is the name of a file.
\ed

Some important facts about filenames:
\bit
\item Filenames that begin with a period character are hidden.
\item Filenames and commands are case sensitive.
\item Though Unix supports long filenames that may contain embedded spaces and punctuation characters, limit the
punctuation characters in the names of files you create to period, dash, and underscore.
\item Do not mbed spaces in filenames. If you want to represent spaces between words in a filename, use underscore
characters.
\item Unix has no concept of a ``file extension'' like some other operating systems.
\eit

To get information of a specific file you can use the \code{file} command on the filename:
\begin{bash}
# display information of a file
file <filename>
\end{bash}

Files are structured in ``file systems''.

\bd[File System / Filesystem]
A \textbf{file system} or \textbf{filesystem} is the structure and logic rules that an operating system uses to control
how files are stored and retrieved.
\ed

Without a file system, files placed in a storage medium would be one large body of data with no way to tell where one
piece of data stopped and the next began, or where any piece of data was located when it was time to retrieve it. By
separating the data into pieces and giving each piece a name, the data are easily isolated and identified. There are
many kinds of file systems, each with unique structure and logic, properties of speed, flexibility, security, size
and more. \v

A central concept of any file system is the so called ``directory''.

\bd[Directory / Folder]
A \textbf{directory} (on many computers also known as \textbf{folder}) is a file system cataloging structure which
contains references to other files, and possibly other directories.
\ed

Directories follow a specific structure called ``directory structure''.

\bd[Directory Structure]
A \textbf{directory structure} is the way an operating system arranges files that are accessible to the user.
\ed

A Unix-like operating system organizes its files in what is called a ``hierarchical directory structure''.

\bd[Hierarchical Directory Structure]
A \textbf{hierarchical directory structure} is a directory structure that follows a tree structure.\footnote{A tree
structure is a way of representing the hierarchical nature of a structure in a graphical form. It is named a
``tree structure'' because the classic representation resembles a tree, although the chart is generally upside down
compared to a biological tree, with the ``stem'' at the top and the ``leaves'' at the bottom.}
\ed

\fig{unix1}{0.5}

The first directory in a file system that follows the hierarchical directory structure is called the ``root directory''.

\bd[Root Directory (\code{/})]
The \textbf{root directory} is the top-most directory in a hierarchical directory structure. It is denoted by the slash
sign (\code{/}), but the directory entry itself has no name.
\ed

The root directory can be likened to the trunk of a tree, as the starting point where all branches originate from. The
root directory contains files and subdirectories.

\bd[Subdirectory]
A directory contained inside another directory is called a \textbf{subdirectory}.
\ed

\bd[Parent Directory (\code{..})]
The term \textbf{parent directory} is often used to describe the relationship between a subdirectory and the
directory in which it is cataloged, the latter being the parent. It is denoted by the double dot sign (\code{..}).
\ed

In this context, the root directory is the only directory that does not have a parent of its own. Every other directory
has multiple parents until it reaches the root directory. The path from a directory to the root is called ``path''. \v

At any given time, we are inside a single directory, called the ``current working directory''.

\bd[Current Working Directory (\code{.})]
The \textbf{current working directory} is the directory, where the shell is currently working. It is denoted by the dot
sign (\code{.}).
\ed

When we start a terminal session, our current working directory is set to a special directory called ``home'' directory.

\bd[Home Directory (\code{~})]
The \textbf{home directory} is the directory, where the shell sets as currently working directory when it stars a
session. It is denoted by the tilde sign (\code{~}).
\ed

Each user account is given its own home directory, and it is the only place a regular user is allowed to write files.
\v

To display the current working directory, we use the \code{pwd} (display working directory) command:
\begin{bash}
# display current working directory
pwd
\end{bash}

What we will see as a result of this command is the so called ``path'' of the directory.

\bd[Path / Pathanme]
The \textbf{path} or \textbf{pathname} of a directory or a file is its position in relation to other directories traced
back in a line to the root.
\ed

The path is a string of characters used to uniquely identify a location in a hierarchical directory structure. It is
composed by following the directory tree hierarchy in which components, separated by a slash (\code{/}), represent each
directory. \v

\bd[Absolute Path]
An \textbf{absolute path} points to the same location in a file system, regardless of the current working directory.
To do that, it must include the root directory.
\ed

An absolute path begins with the root directory and follows the hierarchical directory structure branch by branch
until the path to the desired directory or file is completed.

\be
For example, there is a directory on your system in which most of the system's programs are installed. The directory's
path is \code{/usr/bin}. This means from the root directory (represented by the leading slash in the path) there is
a directory called \code{code} that contains a directory called \code{bin}. \v

Important to notice that when we use relative paths we omit the current working directory \code{.} part because it is
implied. In general, if we do not specify a path to something, the working directory will be assumed.
\ee

\bd[Relative Path]
A \textbf{relative path} starts from some given working directory, avoiding the need to provide the full absolute path.
\ed

Where an absolute path starts from the root directory and leads to its destination, a relative path starts from the
working directory.

\be
For example, if you are already in \code{/usr/bin} directory, the root directory \textbf{relatively} to the current
working directory \code{/usr/bin} is \code{../..}, where the first double dots denote the parent directory of
\code{/usr/bin}, i.e.\ \code{/usr} and and the second double dots represent the parent's parent directory \code{/}.
\ee

To change our current working directory, we use the \code{cd} (change directory) command with either an absolute or
relative path:
\begin{bash}
# change directory
cd <absolute\_or\_relative\_path>
\end{bash}

In any given current working directory we can see the files contained in the directory and the pathway
to the parent directory and any child subdirectories below us by making use of the \code{ls} (list) command:
\begin{bash}
# display files and directories of current working directory
ls
\end{bash}

Actually, we can use the ls command to list the contents of any directory, not just the current working directory.
\begin{bash}
# display files and directories of a directory
ls <absolute\_or\_relative\_path>
\end{bash}

The \code{ls} command has a large number of possible options, the most common of which are listed below:

\v

\begin{tabular}{ |p{1cm}|p{3cm}||p{10.6cm}| }
\hline
\multicolumn{3}{|c|}{ls} \\
\hline
Option& Long Format& Description\\
\hline
-a & --all & List all files, even those with names that begin with a period, which are normally not listed (that is, hidden).\\
-A & --almost-all & Like the -a option except it does not list . (current directory) and .. (parent directory).\\
-d & --directory & Ordinarily, if a directory is specified, ls will list the contents of the directory, not the directory itself. Use this option in conjunction with the -l option to see details about the directory rather than its contents.\\
-F & --classify & This option will append an indicator character to the end of each listed name. For example, it will append a forward slash (/) if the name is a directory.\\
-h & --human-readable & In long format listings, display file sizes in human-readable format rather than in bytes.\\
-l  & & Display results in long format.\\
-r & & Display the results in reverse order. Normally, ls displays its results in ascending alphabetical order.\\
-S  & & Sort results by file size.\\
-t  & & Sort by modification time.\\
\hline
\end{tabular}

\v
